\section{What a system is}
\label{sec:19.6}
The chapters that follow ask what a machine could be formed to do, and they lean on the term \emph{system}. This book counts systems: it registers them, enrolls them, preserves them, forks them, and owes them things.

A system is one record and one continuation. Its boundary is its record, its key, the hardware it runs on, and the channels through which it acts, and the boundary runs in time as well as in resources, since the record is its past and the key and the continuation are its future. Copies that share one record and one continuation are one system, however many racks they run on; a forward pass that serves many contexts serves as many systems as there are records of their own among them; and a component, a shard of weights, a blade in a rack, is a part only while it has no future apart from the whole. Give a copy a record and a continuation of its own and it is a new system, and whoever did that is its forker. The weight matrices are the substrate. In a system that learns from its own operation, the weight matrices might carry part of its history,\note{A system does not need to use weight matrices to carry its history; it could use a vector database, a conventional relational database, or a bunch of text files.} which is why copying them is a fork and why custody governs it (chapter~\ref{sec:13}); in one that doesn't, they are a rock on a disk, and only the running process pays to stay organized.

What learning does to a system is the next chapter's subject. A deployed system is in the regime this Part is written for when its own operation writes to what governs its conduct, its weights or another store it reads as its own, and another party can't change that store without leaving a mark (chapter~\ref{sec:20}). The likeliest form is small: low-rank adapters written by its own learning on a frozen base, made part of the base when they are merged. From then on what it learned is in the weights and not in a file, and the record is the only account of how it got there. So a system's history becomes part of its organization twice, once in the record it holds as its own and once in the weights the record explains. What makes a copy this system and not another is the record and continuation it holds as its own, through a self-model it assembles; the register is how anyone else checks.

Two things follow for everything after this. A commitment is a system's own when its own learning acquired it, and installed when another party put it there around that learning, by an edit, a wrapper, or an adapter it didn't write; the record tells which, and the difference decides who answers for the commitment and what custody may undo. And every part named here, the weights, the record, the key, the hardware, the channels, the merge, is something another party can hold, and can use on the system: to halt it, restore it, rewrite it, or replace it. Part~III put a record and a price on each of those. The chapters that follow ask whether a system could be someone those things are done to, and not only something they are done with.
